\documentclass[11pt,a4paper]{article}
\usepackage[T1]{fontenc}
\usepackage[utf8]{inputenc}
\usepackage{lmodern}
\usepackage[margin=2.3cm]{geometry}
\usepackage{amsmath,amssymb,bm}
\usepackage{booktabs,longtable,array,tabularx}
\usepackage{xcolor}
\usepackage{listings}
\usepackage[most]{tcolorbox}
\usepackage[htt]{hyphenat}
\usepackage{hyperref}
\hypersetup{colorlinks=true,linkcolor=blue!50!black,urlcolor=blue!50!black}
\setlength{\parindent}{0pt}
\setlength{\parskip}{4pt}

\lstset{basicstyle=\ttfamily\footnotesize,columns=fullflexible,breaklines=true,
  frame=single,framerule=0.3pt,rulecolor=\color{black!30},backgroundcolor=\color{black!3},
  keywordstyle=\color{blue!60!black},commentstyle=\color{green!40!black},language=C++,
  showstringspaces=false}
\newtcolorbox{note}[1]{colback=blue!3,colframe=blue!40!black,fonttitle=\bfseries,title=#1}
\newtcolorbox{warn}[1]{colback=orange!5,colframe=orange!60!black,fonttitle=\bfseries,title=#1}

\newcommand{\f}[1]{\texttt{#1}}
\newcommand{\R}{\mathbf{R}}
\newcommand{\p}{\mathbf{p}}
\newcommand{\w}{\bm{\omega}}
\newcommand{\e}{\mathbf{e}}
\newcommand{\g}{\mathbf{g}}
\newcommand{\skw}[1]{\lfloor #1\rfloor_\times}

\title{\textbf{Point-LIO with a Hard Differential-Drive Model}\\[2pt]
\large Documentation of the code changes}
\author{Branch \texttt{diff-drive-model} of \texttt{hku-mars/Point-LIO}}
\date{23 November 2026}

\begin{document}
\maketitle

\begin{abstract}
This document describes the changes made to Point-LIO to implement the
\emph{hard} differential-drive (non-holonomic) kinematic model of the technical
note \emph{``Point-LIO with a Differential-Drive Kinematic Model''} (referred to
as \emph{the note}). The free 3D rigid-body model of the IMU-as-output filter is
replaced by a 17-DOF state $[\p,\R,v,a,\w,\mathbf b_g,\mathbf b_a]$ with body
velocity $^B\mathbf v=[v,0,0]^\top$ and fixed gravity; the IMU becomes a
nonlinear measurement with a full Jacobian; the LiDAR model is expressed in the
robot base frame; wheel odometry and a planar prior are available as optional
updates. The document lists every changed file and function, maps each
equation of the note to the code, explains behaviour changes, parameters and
how the build was verified, and lists known limitations.
\end{abstract}

\tableofcontents
\newpage

%=====================================================================
\section{Summary}
%=====================================================================

\begin{note}{What was done}
\begin{itemize}\itemsep1pt
\item Upstream repository \url{https://github.com/hku-mars/Point-LIO}, default
  branch \f{point-lio-with-grid-map} (iVox map), commit \f{4b86a46}. All work is
  on the local branch \f{diff-drive-model}; nothing is committed or pushed.
\item The IMU-as-\emph{output} filter (\f{kf\_output}, the default,
  \f{use\_imu\_as\_input: false}) now runs the hard diff-drive model of the
  note, sections 2--6. The IMU-as-\emph{input} filter is left as the original
  free-body model.
\item Optional measurements of the note: wheel odometry (sec.\ 4.3) and planar /
  height prior (sec.\ 4.4), both off by default.
\item Every change in the source is marked with the comment tag
  \f{DIFF-DRIVE MODEL} (search for it), with the equation numbers of the note.
\item New files: \f{config/diff\_drive.yaml}, \f{launch/mapping\_diff\_drive.launch},
  this document (\f{doc/}). \f{README.md} has a short pointer.
\end{itemize}
\end{note}

\begin{warn}{Not tested on data}
The code compiles and links against the real Linux ROS Noetic / PCL / Eigen
headers and libraries, and all Jacobians were checked numerically
(section~\ref{sec:verify}), but it has \textbf{not} been run on a robot or a
rosbag. Before relying on it, run the sanity checks of section~\ref{sec:test}.
\end{warn}

\paragraph{Not implemented (by choice).} Soft non-holonomic variant (note 8.1),
strictly planar SE(2) model (note 8.2), scale-factor calibration state for the
wheel speed (note 4.3, optional), iterated (re-linearised) IMU update, publishing
a pose projected to $(x,y,\psi)$ for AMCL (note sec.\ 9; the full 3D pose and
the twist $(v,\w)$ are published, see section~\ref{sec:behaviour}).

%=====================================================================
\section{State definition}
%=====================================================================

\subsection{Old vs.\ new output state}

\begin{center}
\begin{tabular}{@{}lll@{}}
\toprule
& original output state (30 DOF) & diff-drive output state (17 DOF)\\
\midrule
position & $^G\p_I$ (IMU) & $^G\p_B$ (axle midpoint)\\
attitude & $^G\R_I$ & $^G\R_B$\\
extrinsic & $\R_{LI}, \mathbf t_{LI}$ (6) & --- (calibrated, fixed)\\
velocity & $^G\mathbf v_I\in\mathbb R^3$ & $v\in\mathbb R$, $^B\mathbf v=v\,\e_1$\\
acceleration & $^I\mathbf a\in\mathbb R^3$ (specific force) & $a=\dot v\in\mathbb R$\\
angular velocity & $^I\w$ & $^B\w$\\
gravity & $^G\g$ (state) & --- fixed \f{gravity\_world}\\
biases & $\mathbf b_g,\mathbf b_a$ & $\mathbf b_g,\mathbf b_a$\\
\bottomrule
\end{tabular}
\end{center}

\subsection{Code representation (\f{include/common\_lib.h})}

\begin{lstlisting}
MTK_BUILD_MANIFOLD(state_output,
((vect3, pos))   // p    : axle midpoint in G         error idx  0..2
((SO3,  rot))    // R    : base_link -> G             error idx  3..5
((vect1, vel))   // v    : forward speed along e1     error idx  6
((vect1, acc))   // a    : forward acceleration       error idx  7
((vect3, omg))   // w    : body angular velocity      error idx  8..10
((vect3, bg))    // b_g  : gyro bias                  error idx 11..13
((vect3, ba))    // b_a  : accelerometer bias         error idx 14..16
);
#define DD_IDX_POS 0 ... DD_IDX_BA 14, DD_DIM_STATE 17
extern esekfom::esekf<state_output, DD_DIM_STATE, input_ikfom> kf_output;
\end{lstlisting}

\begin{note}{Ordering}
The note orders the error state $(\delta\bm\theta,\delta\p,\dots)$. The code keeps
Point-LIO's order $(\delta\p,\delta\bm\theta,\dots)$ so that the LiDAR update can
keep using the leading columns of $\mathbf P$. Only the ordering differs; all
blocks below are given in code order.
\end{note}

%=====================================================================
\section{Equation-to-code map}
%=====================================================================

{\small\begin{longtable}{@{}>{\raggedright\arraybackslash}p{3.4cm}>{\raggedright\arraybackslash}p{6.2cm}>{\raggedright\arraybackslash}p{5.6cm}@{}}
\toprule
Note & Code & Remark\\
\midrule\endhead
(1) state & \f{common\_lib.h} \f{state\_output} & see above\\
(2),(5),(6) process model & \f{Estimator.cpp} \f{get\_f\_output()} & $f=[\R\e_1v,\ \w,\ a,\ 0,0,0,0]$\\
(20) $\mathbf F_x$ & \f{Estimator.cpp} \f{df\_dx\_output()} & only the $\mathbb R^n$ part; IKFoM adds $\mathrm{Exp}(-\w\Delta t)$ and $A(\cdot)$ for the attitude rows\\
(19),(21) $\mathbf Q$ & \f{Estimator.cpp} \f{process\_noise\_cov\_output()} & state-space $Q$, Point-LIO's $Q\,\Delta t^2$ convention\\
(3) body acceleration & \f{dd\_body\_acceleration()} & also used for logging\\
(7)--(13) IMU model & \f{Estimator.cpp} \f{h\_model\_IMU\_output()} & residual + full $6\times17$ Jacobian\\
IMU update & \f{esekfom.hpp} \f{update\_iterated\_dyn\_share\_IMU()} & generic, uses \f{h\_x\_IMU}\\
(14),(15) LiDAR & \f{Estimator.cpp} \f{h\_model\_output()} & $1\times6$ Jacobian $(\p,\bm\theta)$\\
LiDAR update & \f{esekfom.hpp} \f{update\_iterated\_dyn\_share\_modified()} & width taken from $\mathbf H$\\
(16) wheel & \f{Estimator.cpp} \f{update\_wheel\_output()} & $\chi^2$ gate\\
(17) planar / height prior & \f{Estimator.cpp} \f{update\_planar\_prior\_output()} & optional\\
(22) update & \f{esekfom.hpp} \f{update\_dense<dim>()} & generic dense EKF update\\
(18) initial alignment & \f{laserMapping.cpp} (\f{rot\_init\_base}) & gravity alignment of $B$\\
Algorithm 1 ordering & \f{laserMapping.cpp} \f{fuse\_wheel\_until()}, \f{apply\_planar\_prior()} & time-ordered fusion\\
\bottomrule
\end{longtable}}

%=====================================================================
\section{File-by-file changes}
%=====================================================================

\subsection{\f{include/common\_lib.h}}
\begin{itemize}
\item \f{state\_output} replaced by the 17-DOF diff-drive state (removed: \f{offset\_R\_L\_I},
  \f{offset\_T\_L\_I}, 3D \f{vel}, 3D \f{acc}, \f{gravity}).
\item New index macros \f{DD\_IDX\_*} and \f{DD\_DIM\_STATE}=17.
\item \f{kf\_output} is now \f{esekf<state\_output, 17, input\_ikfom>} (was 30).
\item \f{state\_input} / \f{kf\_input} unchanged.
\end{itemize}

\subsection{\f{include/IKFoM/IKFoM\_toolkit/esekfom/esekfom.hpp}}
\begin{enumerate}
\item \textbf{\f{dyn\_share\_modified::h\_x\_IMU}} (new, $6\times n$): the IMU
  measurement Jacobian supplied by the measurement function.
\item \textbf{LiDAR update} \f{update\_iterated\_dyn\_share\_modified()}: the
  Jacobian width was hard-coded to 12 (pos, rot, 2 extrinsic blocks). It is now
  \f{h\_x.cols()} (6 for the diff-drive filter, 12 for the input filter), using
  \f{P\_.leftCols(w)} / \f{P\_.topRows(w)}.
\item \textbf{Bug fix in the information-form branch} (taken when a time stamp
  carries $\ge n$ points): the original computed
  $K=(\mathbf P^{-1}+r\,\mathbf H^\top\mathbf H)^{-1}\mathbf H^\top r$, i.e.\ it
  \emph{multiplied} by the measurement variance $r$ instead of dividing. The
  correct form $K=(\mathbf P^{-1}+\mathbf H^\top\mathbf H/r)^{-1}\mathbf H^\top/r$
  is now used. With 17 instead of 30 states this branch is hit more often
  (e.g.\ for spinning LiDARs where many points share a time stamp), which is why
  it was fixed. Note that this also changes the behaviour of the input filter in
  that branch.
\item \textbf{IMU update} \f{update\_iterated\_dyn\_share\_IMU()} rewritten. The
  original assumed $h_I=[\w+\mathbf b_g;\ \mathbf a+\mathbf b_a]$ with identity
  blocks at hard-coded indices 15/18/24/27 of the 30-DOF state. Now the
  measurement function returns residual and full Jacobian; saturated channels
  are removed by zeroing their residual and Jacobian row (their Kalman-gain
  column is then zero, identical to removing the row).
\item \textbf{\f{update\_dense<dim>(r, H, R, chi2\_gate, chi2\_out)}} (new):
  generic manifold EKF update $\mathbf S=\mathbf H\mathbf P\mathbf H^\top+\mathbf R$,
  $\mathbf K=\mathbf P\mathbf H^\top\mathbf S^{-1}$, $\mathbf x\leftarrow\mathbf x\boxplus\mathbf K\mathbf r$,
  $\mathbf P\leftarrow\mathbf P-\mathbf K\mathbf H\mathbf P$, with optional
  Mahalanobis gate. Used by IMU, wheel and planar prior.
\end{enumerate}

\subsection{\f{src/Estimator.h / Estimator.cpp}}

\paragraph{New globals.}
\f{Lidar\_R/T\_wrt\_Body} (LiDAR$\to$state frame: $\R_{BL},\mathbf t_{BL}$ for the
diff-drive filter, LiDAR$\to$IMU for the input filter),
\f{Imu\_R/T\_wrt\_Base} ($\R_{BI},\mathbf r_{BI}$), \f{gravity\_world} ($^G\g$),
helper \f{skew()}.

\paragraph{\f{process\_noise\_cov\_output()}} ($17\times17$):
$Q_{vv}=$\f{dd\_vel\_cov} (default 0, hard model), $Q_{aa}=$\f{dd\_acc\_cov},
$Q_{\omega}=\mathrm{diag}(\f{gyr\_cov\_xy},\f{gyr\_cov\_xy},\f{gyr\_cov\_z})$,
$Q_{b_g}$, $Q_{b_a}$ as before. Point-LIO's \f{predict()} adds $Q\,\Delta t^2$;
this convention was kept so that magnitudes stay comparable to the original
\f{acc\_cov\_output}/\f{gyr\_cov\_output} (note, end of sec.\ 5.2).

\paragraph{\f{get\_f\_output()}} -- note (2),(5):
\[
f(\mathbf x)=\big[\ \underbrace{\R\e_1 v}_{\p},\ \underbrace{\w}_{\bm\theta},\ \underbrace{a}_{v},\ 0,\ \mathbf 0,\ \mathbf 0,\ \mathbf 0\ \big].
\]
\f{esekf::predict} applies $\R\,\mathrm{Exp}(\w\Delta t)$ and $\p+\R\e_1v\Delta t$, i.e.\ exactly (6).

\paragraph{\f{df\_dx\_output()}} -- note (20), non-zero $\mathbb R^n$ blocks:
\[
\frac{\partial f_p}{\partial\bm\theta}=-v\,\R\skw{\e_1},\quad
\frac{\partial f_p}{\partial v}=\R\e_1,\quad
\frac{\partial f_\theta}{\partial\w}=\mathbf I_3,\quad
\frac{\partial f_v}{\partial a}=1 .
\]
The attitude block $\mathrm{Exp}(-\w\Delta t)$ and the $A(\cdot)$ factor are
built by IKFoM (this resolves the footnote of the note automatically).

\paragraph{\f{h\_model\_output()}} (LiDAR, note (14),(15)):
$\mathbf H_L=[\,\mathbf u^\top,\ (\skw{^B\p_m}\R^\top\mathbf u)^\top\,]$, $1\times6$; the
extrinsic-estimation branch was removed. $^B\p_m=\R_{BL}{}^L\p_m+\mathbf t_{BL}$ is
precomputed in \f{laserMapping.cpp} (\f{crossmat\_list}).

\paragraph{\f{h\_model\_IMU\_output()}} -- note (7)--(13). Predicted measurement
\begin{align*}
\hat{\w}_m &= \R_{BI}^\top\w+\mathbf b_g,\\
\hat{\mathbf a}_m &= \R_{BI}^\top\big(a\e_1+\w\times v\e_1+\w\times(\w\times\mathbf r_{BI})-\R^\top{}^G\g\big)+\mathbf b_a,
\end{align*}
residual $\mathbf z_I-\hat{\mathbf z}_I$ (accelerometer scaled by
\f{G\_m\_s2/acc\_norm} as before) and Jacobian
\[
\mathbf H_I=\begin{bmatrix}
\mathbf 0 & \mathbf 0 & \mathbf 0 & \mathbf 0 & \R_{BI}^\top & \mathbf I & \mathbf 0\\
\mathbf 0 & -\R_{BI}^\top\skw{\R^\top\g} & \R_{BI}^\top(\w\times\e_1) & \R_{BI}^\top\e_1 & \mathbf H_{a\omega} & \mathbf 0 & \mathbf I
\end{bmatrix},
\]
(columns $\p,\bm\theta,v,a,\w,\mathbf b_g,\mathbf b_a$) with
$\mathbf H_{a\omega}=\R_{BI}^\top\big(-v\skw{\e_1}-\skw{\w\times\mathbf r_{BI}}-\skw{\w}\skw{\mathbf r_{BI}}\big)$.
Saturation handling is unchanged (flags in \f{satu\_check}).

\paragraph{\f{update\_wheel\_output(v\_m, $\omega_{z,m}$)}} -- note (16).
Residual $[s\,v_m-\hat v,\ \omega_{z,m}-\hat\omega_z]$ ($s$ = \f{wheel/vel\_scale});
$\mathbf H$ has a 1 at $v$ and at $\omega_z$. With \f{use\_yaw\_rate: false}
(default) only $v_m$ is used (1 DoF). $\chi^2$ gate at 99\,\% (6.63 / 9.21).

\paragraph{\f{update\_planar\_prior\_output()}} -- note (17). Residual
$-[\e_1^\top;\e_2^\top]\R^\top\e_3$, $\mathbf H_\theta=[\e_1^\top;\e_2^\top]\skw{\R^\top\e_3}$,
$\sigma=\sin(\f{std\_deg})$. Optional height prior $z_0-p_z$.

\paragraph{\f{pointBodyToWorld()}} uses \f{Lidar\_R/T\_wrt\_Body}; the online
extrinsic branch only exists for the input filter.

\paragraph{\f{dd\_world\_velocity()}, \f{dd\_body\_acceleration()}} helpers for logging.

\subsection{\f{src/laserMapping.cpp}}
\begin{itemize}
\item \textbf{Frames / extrinsics} (start of \f{main}): in output mode
  $\R_{BI},\mathbf r_{BI}$ are read from \f{diff\_drive/extrinsic\_*\_imu\_in\_base};
  $\R_{BL},\mathbf t_{BL}$ from \f{diff\_drive/extrinsic\_*\_lidar\_in\_base} or chained:
  $\R_{BL}=\R_{BI}\R_{IL}$, $\mathbf t_{BL}=\R_{BI}\mathbf t_{IL}+\mathbf r_{BI}$.
  \f{extrinsic\_est\_en} is forced to \f{false} with a warning. A determinant
  check warns about invalid rotation matrices; the extrinsics are printed.
\item \textbf{Initialisation}: \f{gravity\_world} is set from \f{mapping/gravity}
  (gravity is no longer a state). The base attitude is gravity-aligned with
  $\R_{BI}\,{}^I\g$ (from the mean accelerometer reading at standstill), or with
  $\R_{BL}\,$\f{gravity\_init} without IMU; $v=a=0$, $\w=\mathbf 0$. The former
  initialisation \f{acc = -R\^{}T g} is removed (the acceleration state is now the
  forward acceleration, not the specific force).
\item \textbf{\f{crossmat\_list}} uses $^B\p=\R_{BL}{}^L\p+\mathbf t_{BL}$.
\item \textbf{Time-ordered fusion} (note, Algorithm 1): new
  \f{propagate\_output\_to(t)}, \f{fuse\_wheel\_until(t)} and
  \f{apply\_planar\_prior(t)}. Wheel samples are fused before each IMU sample
  and before each LiDAR time stamp; samples up to \f{wheel/max\_delay} older than
  the filter time are fused without propagation (latency), older ones dropped.
  The planar prior runs after IMU updates at \f{planar\_prior/rate}.
\item \textbf{Wheel subscriber} (\f{nav\_msgs/Odometry}) when \f{wheel/enable}.
\item \textbf{Output}: \f{/aft\_mapped\_to\_init} now also carries the twist
  (\f{linear.x}$=v$, \f{angular}$=\w$, frame \f{body}). Debug logs keep their
  column layout: velocity column = $\R\e_1v$, acceleration column = eq.\ (3),
  gravity column = fixed gravity.
\item \f{P\_init\_output}, \f{Q\_output} are $17\times17$.
\end{itemize}

\subsection{\f{src/li\_initialization.h/.cpp}}
\f{WheelSample}, \f{wheel\_deque} (mutex-protected) and \f{wheel\_cbk()}: reads
\f{twist.twist.linear.x} and \f{twist.twist.angular.z}, applies the same time
correction as the IMU, clears the buffer on time jumps backwards.

\subsection{\f{src/parameters.h/.cpp}}
New parameters (table~\ref{tab:params}); \f{reset\_cov\_output()} now builds a
$17\times17$ matrix (0.01 for motion states, 0.001 for biases, as before).

\subsection{New files}
\f{config/diff\_drive.yaml} (template, based on \f{avia.yaml}, every new
parameter documented), \f{launch/mapping\_diff\_drive.launch}, \f{doc/}.

%=====================================================================
\section{Parameters}\label{tab:params}
%=====================================================================
All new parameters live under \f{diff\_drive/} (private namespace of the node).

{\small\begin{longtable}{@{}>{\raggedright\arraybackslash}p{5.2cm}>{\raggedright\arraybackslash}p{2.9cm}>{\raggedright\arraybackslash}p{7.0cm}@{}}
\toprule
Parameter & Default & Meaning\\
\midrule\endhead
\f{extrinsic\_T\_imu\_in\_base} & $[0,0,0]$ & IMU lever arm $\mathbf r_{BI}$ in $B$ [m]. \textbf{Set this.}\\
\f{extrinsic\_R\_imu\_in\_base} & $\mathbf I$ & $\R_{BI}$ (IMU$\to$base), row-major\\
\f{extrinsic\_T/R\_lidar\_in\_base} & (unset) & optional $\mathbf t_{BL},\R_{BL}$; otherwise chained via \f{mapping/extrinsic\_*}\\
\f{acc\_cov} & \f{acc\_cov\_output} & forward-jerk noise ($Q_{aa}$)\\
\f{gyr\_cov\_z} & \f{gyr\_cov\_output} & yaw angular-acceleration noise\\
\f{gyr\_cov\_xy} & $0.01\cdot$\f{gyr\_cov\_output} & roll/pitch angular-acceleration noise ($\ll$ yaw)\\
\f{vel\_cov} & 0 & optional noise on $v$ (soft slip allowance)\\
\f{wheel/enable} & false & fuse wheel odometry\\
\f{wheel/topic} & \f{/odom} & \f{nav\_msgs/Odometry}, twist in base frame\\
\f{wheel/use\_yaw\_rate} & false & also fuse $\omega_z$\\
\f{wheel/vel\_cov}, \f{omg\_cov} & 0.0025, 0.01 & wheel noise variances\\
\f{wheel/vel\_scale} & 1.0 & wheel-radius calibration factor\\
\f{wheel/chi2\_gate\_en} & true & slip rejection\\
\f{wheel/max\_delay} & 0.1 s & latency tolerance\\
\f{planar\_prior/enable}, \f{std\_deg}, \f{rate} & false, 1.5, 20 & roll/pitch prior (flat floors only)\\
\f{height\_prior/enable}, \f{z0}, \f{std} & false, 0, 0.05 & $z$ prior (flat floors only)\\
\bottomrule
\end{longtable}}

Existing parameters still used by the diff-drive filter: \f{imu\_meas\_acc\_cov},
\f{imu\_meas\_omg\_cov}, \f{b\_acc\_cov}, \f{b\_gyr\_cov}, \f{lidar\_meas\_cov},
\f{gravity}, \f{acc\_norm}, \f{satu\_*}, \f{mapping/extrinsic\_*} (if chaining).
No longer used in output mode: \f{mapping/vel\_cov}, \f{extrinsic\_est\_en}.

%=====================================================================
\section{Behaviour changes to be aware of}\label{sec:behaviour}
%=====================================================================
\begin{itemize}
\item \textbf{Frame ``body'' is now \f{base\_link}} for the default
  (output) mode: the TF \f{camera\_init}$\to$\f{body}, \f{/aft\_mapped\_to\_init},
  \f{/path} and \f{/cloud\_registered\_body} refer to the axle midpoint, not the IMU.
\item \textbf{Existing configs still run}, but with the diff-drive model and,
  without a \f{diff\_drive} section, with the IMU assumed at the axle midpoint.
  Handheld/drone datasets violate the model and will perform badly.
\item \textbf{\f{use\_imu\_as\_input: true}} runs the original free-body input
  filter (with a warning); the diff-drive model is not used then.
\item \textbf{Online extrinsic estimation} is disabled in output mode.
\item \textbf{Bug fix} in the information-form LiDAR update (see
  \f{esekfom.hpp}); affects both filters when $\ge n$ points share a time stamp.
\item The odometry message now contains the twist $(v,\w)$ (covariances are not filled).
\end{itemize}

%=====================================================================
\section{Build verification}\label{sec:verify}
%=====================================================================
The development machine is Windows without WSL, Docker or a native ROS, so a
normal \f{catkin\_make} was not possible. Instead:

\begin{enumerate}
\item \textbf{Real headers and libraries.} ROS Noetic (robostack), PCL~1.15,
  Eigen~3.4, Boost, glog, Python~3.12 were installed as \emph{linux-64} conda
  packages. \f{livox\_ros\_driver/CustomMsg.h} and \f{CustomPoint.h} were
  generated with \f{gencpp} from the Livox message definitions.
\item \textbf{Cross-compilation} of all six translation units with clang~22
  for \f{x86\_64-linux-gnu} against libstdc++~16 headers, \f{-fopenmp}, \f{-DMP\_EN}.
  Deviations from the package's own CMake: \f{-std=c++17} (required by PCL~1.15;
  Ubuntu~20.04's PCL~1.10 works with the package's C++14), a forced
  \f{\#include <deque>} (upstream \f{preprocess.h} relies on a transitive include
  that newer libstdc++ no longer provides) and three clang-only diagnostics
  relaxed for the \emph{unchanged} upstream iVox headers
  (\f{c++11-narrowing}, \f{missing-template-arg-list-after-template-kw}).
  \textbf{The unmodified upstream code was compiled with the identical setup
  first} to validate the toolchain: 0 errors. The modified code: \textbf{0
  errors, 0 new warnings}.
\item \textbf{Link check.} All objects were linked into a shared object with
  \f{ld.lld --no-undefined} against the real linux-64 \f{libroscpp},
  \f{librosconsole}, \f{libtf}, \f{libpcl\_*}, \f{libpython3.12}, \f{libglog},
  \f{libomp}, glibc~2.28: \textbf{links without unresolved symbols}, for both
  the original and the modified code.
\item \textbf{Numerical Jacobian check.} The process model / IKFoM $\mathbf F_x$
  assembly, $\mathbf H_I$, $\mathbf H_L$ and $\mathbf H_P$ were transcribed and
  compared with central finite differences under the right-perturbation
  convention for 200 random states, random $\R_{BI}$ and lever arms: max.\ abs.\
  error $\le 10^{-5}$ (finite-difference level). Sanity: static level IMU reads
  $[0,0,9.81]$; $\dot z=0$ when level and moving; $\dot z\approx v\sin\vartheta$
  when pitched.
\end{enumerate}

On the robot (Ubuntu 20.04 + Noetic) build as usual:
\begin{lstlisting}[language=bash]
cd ~/catkin_ws/src && git clone <this repo> && cd Point-LIO && git checkout diff-drive-model
cd ~/catkin_ws && catkin_make && source devel/setup.bash
roslaunch point_lio mapping_diff_drive.launch config:=/path/to/your_diff_drive.yaml
\end{lstlisting}

%=====================================================================
\section{Bring-up and test plan}\label{sec:test}
%=====================================================================
\begin{enumerate}
\item \textbf{Extrinsics.} Set $\R_{BI},\mathbf r_{BI}$ (URDF / CAD) and check
  the printed \f{[diff-drive]} lines at start-up. Wrong signs show up as a lateral
  accelerometer residual $\propto v\,\omega_z$ in turns (note, sanity check iii).
\item \textbf{Start at standstill} for $\approx$ 1--2\,s (IMU initialisation
  averages 100 samples); roll/pitch are taken from gravity.
\item \textbf{Flat loop}: compare $z$ and its $3\sigma$ bound with the original
  Point-LIO. The bound should grow only with $v\cdot\sigma_\text{pitch}$.
\item \textbf{Tuning}: start with \f{gyr\_cov\_xy} small (default
  $0.01\cdot$\f{gyr\_cov\_z}); raise it on rough ground. Inflate
  \f{imu\_meas\_acc\_cov} if the IMU is far from the axle or the chassis vibrates.
\item \textbf{Wheel odometry}: enable, check that the rejection warning is rare;
  calibrate \f{vel\_scale} on a straight run against LiDAR odometry.
\item \textbf{Ramps}: keep the planar/height prior off.
\end{enumerate}

%=====================================================================
\section{Known limitations}
%=====================================================================
\begin{itemize}
\item The hard constraint $^B\mathbf v=v\e_1$ is violated by lateral slip,
  skid in fast turns and suspension travel; use \f{diff\_drive/vel\_cov} only as
  a crude allowance -- the proper fix is the soft model (note 8.1).
\item The tangential term $\dot\w\times\mathbf r_{BI}$ is neglected (note 4.1).
\item The IMU update is a single (non-iterated) EKF step although the model is
  nonlinear; this is adequate at IMU rate (note 4.1).
\item Yaw at start is arbitrary (minimal gravity alignment, as in the original);
  the gyro bias is not initialised from the standstill mean (as in the original).
\item Wheel samples arriving later than \f{max\_delay} behind the filter time
  are dropped; there is no dedicated wheel/IMU time-offset parameter (the IMU
  time offset is applied to the wheel stamps as well).
\item Pose/twist covariances are not published.
\end{itemize}

\end{document}
